\begin{enumerate}[label=\textbf{1.1.\arabic*.},ref=1.1.\arabic*]
\setlength{\itemsep}{0.35em}
\item \label{formula-1-1-1} \textbf{Quadratic equations.} For $ax^2+bx+c=0$,
\begin{equation*}\tag{1.1.1.1}
\Delta=b^2-4ac,\qquad x=\frac{-b\pm\sqrt{\Delta}}{2a},\qquad
\alpha+\beta=-\frac ba,\quad \alpha\beta=\frac ca.
\end{equation*}
Equal real roots occur when $\Delta=0$. For $y=ax^2+bx+c$, the vertex has $x$-coordinate $-b/(2a)$.
\item \label{formula-1-1-2} \textbf{Tangents to a circle.} A radius to the point of contact is perpendicular to the tangent. Tangents from the same external point are equal, and for radius $r$,
\begin{equation*}\tag{1.1.2.1}
PT^2=OP^2-r^2.
\end{equation*}
\item \label{formula-1-1-3} \textbf{Proportionality and similarity.} If $AB\parallel EF$ in $\triangle DEF$, then
\begin{equation*}\tag{1.1.3.1}
\frac{DA}{DE}=\frac{DB}{DF}=\frac{AB}{EF},\qquad
\frac{DA}{AE}=\frac{DB}{BF}.
\end{equation*}
For similar triangles, corresponding side ratios are equal and their area ratio is the square of that ratio.
\item \label{formula-1-1-4} \textbf{Coordinate geometry.} For $A(x_1,y_1)$ and $B(x_2,y_2)$,
\begin{equation*}\tag{1.1.4.1}
AB=\sqrt{(x_2-x_1)^2+(y_2-y_1)^2},\qquad
M=\left(\frac{x_1+x_2}{2},\frac{y_1+y_2}{2}\right).
\end{equation*}
If $AP:PB=m:n$, then $P=\left(\frac{nx_1+mx_2}{m+n},\frac{ny_1+my_2}{m+n}\right)$.
\item \label{formula-1-1-5} \textbf{Slope and parallel lines.} The slope through two points is $m=(y_2-y_1)/(x_2-x_1)$. Non-vertical parallel lines have equal slopes.
\item \label{formula-1-1-6} \textbf{Circle sectors and chords.} For a central angle $\theta$ in degrees,
\begin{equation*}\tag{1.1.6.1}
\text{arc length}=\frac{\theta}{360^\circ}2\pi r,\qquad
\text{sector area}=\frac{\theta}{360^\circ}\pi r^2.
\end{equation*}
A segment's area is the sector area minus the area of its triangle.
\item \label{formula-1-1-7} \textbf{Right-triangle trigonometry.} $\sin\theta=\frac{\text{opposite}}{\text{hypotenuse}}$, $\cos\theta=\frac{\text{adjacent}}{\text{hypotenuse}}$, and $\tan\theta=\frac{\text{opposite}}{\text{adjacent}}$.
\item \label{formula-1-1-8} \textbf{Mensuration.}
\begin{equation*}\tag{1.1.8.1}
V_{\rm cone}=\frac13\pi r^2h,\qquad
V_{\rm sphere}=\frac43\pi r^3,\qquad
S_{\rm sphere}=4\pi r^2,\qquad
S_{\rm cylinder,curved}=2\pi rh.
\end{equation*}
\item \label{formula-1-1-9} \textbf{Probability.} For equally likely outcomes, $P(E)=\frac{n(E)}{n(S)}$; $P(E^c)=1-P(E)$.
\item \label{formula-1-1-10} \textbf{HCF and LCM.} For positive integers $a,b$, $\operatorname{HCF}(a,b)\operatorname{LCM}(a,b)=ab$; the Euclidean algorithm finds the HCF by repeated division.
\end{enumerate}
